\section*{References（参考文献）}
\addcontentsline{toc}{section}{References（参考文献）}

\subsection*{聖書（旧約聖書）}
\noindent 日本聖書協会

\begin{quote}
本研究では、旧約聖書における「詩篇」や賛美文化を、人類が共同体の中で歌い、祈り、感情を共有してきた文化的背景として参照した。

また、人類が古くから「死」や「祝福」「季節の節目」などを、単なる出来事ではなく、歌・祈り・儀式として共同体の中で共有してきた背景にも着目している。

本研究では、そのような行為は宗教的形式を超えて、人間が本来的に持っている「共に時間を過ごし、感情を分かち合おうとする性質」の一つである可能性を考察する上で参考とした。
\end{quote}

\subsection*{聖書（新約聖書）}
\noindent 日本聖書協会

\begin{quote}
特に「ルカによる福音書」におけるキリスト降誕記述を中心に、クリスマス文化の基盤として参照した。本研究では、宗教そのものではなく、「共に祝う」「共に歌う」「希望を共有する」という文化構造に着目している。
\end{quote}

\subsection*{Bowler, Gerry.}
\noindent Santa Claus: A Biography.\\McClelland \& Stewart, 2005.

\begin{quote}
サンタクロース文化の歴史的背景、聖ニコラス伝承、金貨と靴下に関する逸話などを参照した。本研究では、「文化背景を知ることによって参加者の愛着や理解が深まる構造」を考察する上で参考とした。
\end{quote}

\subsection*{Pine, B. Joseph \& Gilmore, James H.}
\noindent The Experience Economy.\\Harvard Business School Press, 1999.

\begin{quote}
「体験経済（Experience Economy）」理論を参照した。本研究では、観光を単なるサービス消費ではなく、「共同参加によって生成される体験」として捉える視点において関連性を持つ。
\end{quote}

\subsection*{Richards, Greg.}
\noindent Creative Tourism and Local Development.\\Creative Tourism: A Global Conversation, 2007.

\begin{quote}
参加型・創造型観光研究を参照した。本研究では、観光客が単なる鑑賞者ではなく、地域文化へ主体的に参加する構造について考察する際の参考とした。
\end{quote}

\subsection*{Putnam, Robert D.}
\noindent Bowling Alone: The Collapse and Revival of American Community.\\Simon \& Schuster, 2000.

\begin{quote}
現代社会における共同体参加や地域接続の減少について論じた研究。本研究では、「共に歌う」という行為によって地域住民・観光客・ホテル関係者が再接続される構造を考察する上で参照した。
\end{quote}

\subsection*{Oldenburg, Ray.}
\noindent The Great Good Place.\\Marlowe \& Company, 1999.

\begin{quote}
「Third Place（第三の場所）」理論を参照した。本研究では、太陽の教会が、家庭や職場とは異なる「人々が自然につながる場所」として機能した点との関連性を持つ。
\end{quote}

\subsection*{Manzini, Ezio.}
\noindent Design, When Everybody Designs.\\MIT Press, 2015.

\begin{quote}
市民参加型デザインおよび社会設計思想を参照した。本研究では、地域住民・観光客・ホテルが共に空間や体験を形成していく「共同設計構造」を考察する際の参考とした。
\end{quote}
